\section{Entropic Regularization: Convergence}
\label{sec-sinkhorn-advanced}
\label{sec-entropic-convergence}
\label{sec-convergence-dual}
The chapter revisits Sinkhorn convergence through several complementary lenses. Bregman projections explain the alternating-projection geometry, while Fortet's order argument gives qualitative fixed-point convergence.

\subsection{Sinkhorn Convergence: Bregman Point of View}
\label{sec-convergence-init}
\paragraph{\texorpdfstring{Alternating $\KL$ projections.}{Alternating KL projections.}}
\begin{defn}[Bregman divergence]\label{def-bregman-divergence}
Let $\Omega\subset\RR^{n\times m}$ be convex with nonempty interior, and let $\Phi:\RR^{n\times m}\to\RR\cup\{+\infty\}$ be a proper lower-semicontinuous Legendre function with $\operatorname{dom}(\Phi)=\Omega$. Thus $\Phi=+\infty$ outside $\Omega$, while $\Phi$ is differentiable and strictly convex on $\operatorname{int}(\Omega)$. For $\P,\Q\in\operatorname{int}(\Omega)$, the Bregman divergence generated by $\Phi$ is
\(B_\Phi(\P|\Q)\eqdef \Phi(\P)-\Phi(\Q)-\dotp{\nabla\Phi(\Q)}{\P-\Q}.\)
\end{defn}
\begin{defn}[Bregman projection]\label{def-bregman-projection}
For a nonempty convex set $\Cc$, define
\(\Proj_{\Cc}^{B_\Phi}(\Q) \eqdef \argmin_{\P\in\Cc} B_\Phi(\P|\Q),\)
which is set-valued if the minimizer is not unique.
\end{defn}
\paragraph{Linear tilts and Gibbs references.}
\begin{prop}[Linear tilts of Bregman penalties]\label{prop-bregman-linear-tilt}
Let $\Phi$ be differentiable and strictly convex, and let $B_\Phi$ be its Bregman divergence. Fix a reference point $\Q$ in the interior of the domain. Assume that there exists $\Q^\C$ such that
\(\nabla\Phi(\Q^\C)=\nabla\Phi(\Q)-\C/\epsilon.\)
Then, for all $\P$ in the domain, \(\dotp{\P}{\C}+\epsilon B_\Phi(\P|\Q) = \epsilon B_\Phi(\P|\Q^\C)+\text{\upshape cst},\)
where the constant does not depend on $\P$.
\end{prop}
\begin{proof}
Subtract the two Bregman divergences: \(B_\Phi(\P|\Q^\C)-B_\Phi(\P|\Q) = \dotp{\nabla\Phi(\Q)-\nabla\Phi(\Q^\C)}{\P} +\text{cst}.\)
Using $\nabla\Phi(\Q)-\nabla\Phi(\Q^\C)=\C/\epsilon$ and multiplying by $\epsilon$ gives the claim.
\end{proof}
For the negative entropy $\Phi(\P)=\sum_{i,j}\P_{i,j}\log\P_{i,j}$, one has $B_\Phi=\KLD$. Taking $\Q=\a\otimes\b$ gives the tilted reference

\(\K_{\a,\b}^\epsilon \eqdef (\a\otimes\b)\odot e^{-\C/\epsilon}.\)
\(\dotp{\P}{\C}+\epsilon\KLD(\P|\a\otimes\b) = \epsilon\KLD(\P|\K_{\a,\b}^\epsilon)+\text{\upshape cst}.\)
Thus the unique solution $\P_\epsilon$ of~\eqref{eq-regularized-discr} is the KL projection of the tilted Gibbs reference onto $\CouplingsD(\a,\b)$:

\eql{\label{eq-kl-proj}
\P_\epsilon = \Proj_{\CouplingsD(\a,\b)}^\KLD(\K_{\a,\b}^\epsilon) \eqdef \uargmin{\P \in \CouplingsD(\a,\b)} \KLD(\P|\K_{\a,\b}^\epsilon).
}
\paragraph{Cyclic projection convergence.}
Given two closed convex constraint sets $\Cc_1$ and $\Cc_2$, the cyclic method projects first onto $\Cc_1$, then onto $\Cc_2$, and repeats. Full iterates satisfy the second constraint and half-iterates satisfy the first; both constraints are attained together only in the limit.

\begin{prop}[Convergence of cyclic Bregman projections]\label{prop-cyclic-kl-affine}
Let $\Phi$ be a Legendre generator on a finite-dimensional convex domain, and let $\Cc_1,\Cc_2$ be closed convex sets whose intersection meets the interior of the domain. Generate $(\P^{(\ell-1/2)},\P^{(\ell)})$ by Algorithm~\ref{alg:cyclic-bregman-projections} from an interior point $\P^{(0)}$. Assume that the projections are well-defined and that the full half-step sequence remains in a compact subset of the interior of the domain. Then there exists $\bar\P\in\Cc_1\cap\Cc_2$ such that $\P^{(\ell)}\to\bar\P$ and $\P^{(\ell-1/2)}\to\bar\P$.
If, in addition, $\Cc_1$ and $\Cc_2$ are affine, then
\(\bar\P = \Proj_{\Cc_1\cap\Cc_2}^{B_\Phi}(\P^{(0)}).\)
\end{prop}
\begin{proof}
For three interior points, the Bregman three-point formula reads \(B_\Phi(\Q|\P) = B_\Phi(\Q|\P^+) + B_\Phi(\P^+|\P) + \dotp{\nabla\Phi(\P^+)-\nabla\Phi(\P)}{\Q-\P^+}.\)
If $\P^+=\Proj_\Cc^{B_\Phi}(\P)$ and $\Cc$ is closed and convex, first-order optimality gives \(\dotp{\nabla\Phi(\P^+)-\nabla\Phi(\P)}{\Q-\P^+}\geq0 \qquad \forall \Q\in\Cc.\)
\(B_\Phi(\Q|\P) \geq B_\Phi(\Q|\P^+)+B_\Phi(\P^+|\P) \qquad\forall \Q\in\Cc.\)

Let $(Z^{(r)})_r$ be the half-step sequence, so that $Z^{(2\ell)}=\P^{(\ell)}$ and $Z^{(2\ell+1)}=\P^{(\ell+1/2)}$. Fix $\Q\in\Cc_1\cap\Cc_2$. Applying the inequality at each half-step gives
\(B_\Phi(\Q|Z^{(r)})-B_\Phi(\Q|Z^{(r+1)}) \geq B_\Phi(Z^{(r+1)}|Z^{(r)})\geq0.\)
Thus $B_\Phi(\Q|Z^{(r)})$ decreases and $\sum_r B_\Phi(Z^{(r+1)}|Z^{(r)})<+\infty$. The projection drops tend to zero; compactness and strict convexity imply $\norm{Z^{(r+1)}-Z^{(r)}}\to0$. Every cluster point of either subsequence consequently belongs to both closed sets. Let $\bar\P\in\Cc_1\cap\Cc_2$ be one such cluster point. The sequence $B_\Phi(\bar\P|Z^{(r)})$ is nonincreasing and tends to zero along a convergent subsequence, hence tends to zero altogether. Compactness and strict convexity then give $Z^{(r)}\to\bar\P$, proving the first claim.

Suppose now that $\Cc_1$ and $\Cc_2$ are affine. The first-order inequality above is then an equality because both signs of every feasible direction are admissible. Moreover, each dual displacement
\(\nabla\Phi(Z^{(r+1)})-\nabla\Phi(Z^{(r)})\)
belongs to the normal space of the affine set used at that half-step. Telescoping and passing to the limit gives \(\nabla\Phi(\bar \P)-\nabla\Phi(\P^{(0)}) \in N_{\Cc_1}+N_{\Cc_2} = N_{\Cc_1\cap\Cc_2},\)
where the last equality uses affinity and the nonempty intersection. This is precisely the first-order optimality condition for minimizing $\Q\mapsto B_\Phi(\Q|\P^{(0)})$ over $\Cc_1\cap\Cc_2$, proving the displayed projection identity.
\end{proof}
\paragraph{Row and column scalings.}
\eq{\label{eq-affine-marginal-sets}
\Cc^1_\a \eqdef \enscond{\P\in\RR^{n\times m}}{\P\ones_m=\a}
\qandq
\Cc^2_\b \eqdef \enscond{\P\in\RR^{n\times m}}{\transp{\P}\ones_n=\b}.
}
\(\CouplingsD(\a,\b) = \Cc^1_\a\cap\Cc^2_\b\cap\RR_+^{n\times m}.\)
For negative entropy, however, the effective domain of the Legendre generator is $\Omega=\RR_+^{n\times m}$ and $\Phi=+\infty$ outside $\Omega$. Thus the two half-steps of Algorithm~\ref{alg:cyclic-bregman-projections} specialize to

\eql{\label{eq-kl-sinkh-proj}
\itt{\P} \eqdef \Proj_{\Cc^1_\a}^{\KLD}(\it{\P})
\qandq
\ittt{\P} \eqdef \Proj_{\Cc^2_\b}^{\KLD}(\itt{\P}).
}
\begin{prop}[KL projections are scalings]
One has
\(\Proj_{\Cc^1_\a}^{\KLD}(\P) = \diag\pa{\frac{\a}{\P \ones_m}} \P \qandq \Proj_{\Cc^2_\b}^{\KLD}(\P) = \P \diag\pa{\frac{\b}{\P^\top \ones_n}}.\)
\end{prop}
\begin{proof}
Consider the problem along each row or column vector to impose a fixed sum $s \in \RR_+$
\(\umin{p} \enscond{\KL(p|q)}{\dotp{p}{\ones}=s}\).
The Lagrange multiplier equation for this problem reads \(\log(p/q)+\la \ones=0 \qarrq p = u q \qwhereq u = e^{-\la}>0.\)
The constraint $\dotp{p}{\ones} = s$ is equivalent to $\dotp{uq}{\ones}=s$, i.e. $u=s/\sum_i q_i$, which gives the desired scaling formula
$p=s q/\sum_i q_i$.
\end{proof}
\eq{\label{eq-sink-matrix}\P^{(2\ell)} \eqdef \diag(\it{\uD}) \K \diag(\it{\vD}),}
\begin{align*}
\P^{(2\ell+1)} &\eqdef \diag(\itt{\uD}) \K \diag(\it{\vD}) \\
\qandq
\P^{(2\ell+2)} &\eqdef \diag(\itt{\uD}) \K \diag(\itt{\vD})
\end{align*}
\paragraph{Other divergences.}
\eql{\label{eq-positive-marginal-sets}
\Cc^1_{\a,+}
\eqdef
\Cc^1_\a\cap\RR_+^{n\times m}
\qandq
\Cc^2_{\b,+}
\eqdef
\Cc^2_\b\cap\RR_+^{n\times m},
}
\(\big[\Proj_{\Cc^1_{\a,+}}^{B_\Phi}(\Q)\big]_{i,j} = (\Q_{i,j}-\tau_i)_+, \qquad \sum_j(\Q_{i,j}-\tau_i)_+=a_i,\)
For the quadratic generator and the product reference $\xi=\a\otimes\b$, the Bregman penalty is

\(B_\Phi(\P|\xi) = \frac12\sum_{i,j}(\P_{i,j}-a_i b_j)^2.\)
\(\Divergm_\phi(\P|\a\otimes\b) = \frac12\sum_{i,j}\frac{(\P_{i,j}-a_i b_j)^2}{a_i b_j}.\)
\subsection{Sinkhorn Convergence: Monotone Point of View}\label{sec-sinkhorn-monotone}
\paragraph{Variation seminorm and topical maps.}
\begin{defn}[Variation seminorm]\label{def-variation-seminorm}
For a bounded real-valued function $h$, its variation seminorm is \(\norm{h}_V\eqdef \sup h-\inf h.\)
It vanishes exactly on constant functions and therefore induces a norm after quotienting by additive constants.
\end{defn}
\begin{defn}[Topical map]\label{def-topical-map}
Let $E$ be an ordered vector space of functions containing the constants. A map $\mathcal T:E\to E$ is \emph{topical} if
\(f\leq g\Longrightarrow\mathcal T(f)\leq\mathcal T(g), \qquad \mathcal T(f+s)=\mathcal T(f)+s \quad(s\in\RR).\)
\end{defn}
\begin{prop}[Topical maps are variation-nonexpansive]\label{prop-topical-variation-nonexpansive}
Let $E$ be a vector space of bounded real-valued functions, ordered pointwise. Every topical map $\mathcal T:E\to E$ satisfies \(\norm{\mathcal T(f)-\mathcal T(g)}_V \leq \norm{f-g}_V.\)
The same conclusion holds if $\mathcal T$ is order reversing and additively anti-homogeneous, meaning $\mathcal T(f+s)=\mathcal T(f)-s$.
\end{prop}
\begin{proof}
Set $a=\inf(f-g)$ and $b=\sup(f-g)$, so that $g+a\leq f\leq g+b$. If $\mathcal T$ is topical, then
\(\mathcal T(g)+a\leq\mathcal T(f)\leq\mathcal T(g)+b.\)
Hence the oscillation of $\mathcal T(f)-\mathcal T(g)$ is at most $b-a=\norm{f-g}_V$. In the order-reversing, additively anti-homogeneous case, the corresponding bounds are $\mathcal T(g)-b\leq\mathcal T(f)\leq\mathcal T(g)-a$ and give the same conclusion.
\end{proof}
\paragraph{Generalized Sinkhorn maps.}
Consider the $\phi$-divergence regularized transport problem of Section~\ref{sec-sinkhorn-other-regularizers}. Its one-sided block updates, expressed using the Legendre transform $\phi^*$ defined in~\eqref{eq-legendre}, are the generalized soft $c$-transforms of~\eqref{eq-phi-soft-c-transform}; their alternating dual-maximization iteration is~\eqref{eq-generalized-soft-c-alternate-maximization}.

\eql{\label{eq-generalized-sinkhorn-map}
\mathcal A_\phi(f)
\eqdef
\big(f^{c,\epsilon,\phi}\big)^{\bar c,\epsilon,\phi}.
}
\begin{prop}[Generalized Sinkhorn maps are topical]\label{prop-phi-double-soft-transform-monotone}
Assume that the scalar minimizer sets in~\eqref{eq-phi-soft-c-transform} are nonempty compact intervals, and select either their smallest elements everywhere or their largest elements everywhere. Each selected one-sided transform is order reversing and additively anti-homogeneous:
\[
g\leq g'
\Longrightarrow
g^{\bar c,\epsilon,\phi}\geq {g'}^{\bar c,\epsilon,\phi},
\qquad
(g+s)^{\bar c,\epsilon,\phi}
=
g^{\bar c,\epsilon,\phi}-s,
\]
and likewise for $f\mapsto f^{c,\epsilon,\phi}$.
\(\norm{\mathcal A_\phi(f)-\mathcal A_\phi(h)}_V \leq \norm{f-h}_V.\)
No selection convention is needed when the scalar minimizers are unique.
\end{prop}
\begin{proof}
Since $\phi$ is extended by $+\infty$ on $(-\infty,0)$, its Legendre transform $\phi^*$ from~\eqref{eq-legendre} is convex and nondecreasing. For fixed $x$, let $H_g^x(u)$ denote the scalar objective in the first line of~\eqref{eq-phi-soft-c-transform}. If $g\leq g'$, then $u\mapsto H_{g'}^x(u)-H_g^x(u)$ is nondecreasing: this follows by integrating the fact that $s\mapsto\phi^*(s+\delta)-\phi^*(s)$ is nondecreasing for every $\delta\geq0$.

Let $I=\argmin H_g^x$ and $I'=\argmin H_{g'}^x$. If $a=\min I$ and $a'=\min I'$ satisfied $a'>a$, the optimality inequalities $H_g^x(a)\leq H_g^x(a')$ and $H_{g'}^x(a')\leq H_{g'}^x(a)$ would contradict the preceding monotonicity; in the equality case, $a$ would also belong to $I'$, contradicting the definition of $a'$. Hence $\min I'\leq\min I$. The same argument with the maximal minimizers gives $\max I'\leq\max I$. Finally, $H_{g+s}^x(u)=H_g^x(u+s)+s$, so either extremal selection is additively anti-homogeneous. The second one-sided transform is treated symmetrically. Their composition is topical, and Proposition~\ref{prop-topical-variation-nonexpansive} gives the final estimate.
\end{proof}
\paragraph{Monotone convergence for generalized Sinkhorn.}
\begin{prop}[Monotone convergence of generalized Sinkhorn cycles]\label{prop-fortet-monotone}
Let $\Xx$ and $\Yy$ be compact metric spaces and $c\in\Cc(\Xx\times\Yy)$. Suppose that the assumptions of Proposition~\ref{prop-phi-double-soft-transform-monotone} hold, that the selected transforms are finite valued, and that $\mathcal A_\phi$ has a fixed point $f^\star\in\Cc(\Xx)$.
If $f^{(0)}\in\Cc(\Xx)$ is a subsolution, $f^{(0)}\leq\mathcal A_\phi(f^{(0)})$, then the iterates $f^{(\ell+1)}=\mathcal A_\phi(f^{(\ell)})$ converge uniformly to a fixed point of $\mathcal A_\phi$. The analogous conclusion holds for a supersolution, with a decreasing sequence. If the fixed point is unique modulo additive constants, the corresponding potential classes converge to this unique class.
\end{prop}
\begin{proof}
Because $\mathcal A_\phi$ is additively homogeneous, shifting a subsolution preserves the subsolution inequality. After such a shift, compactness of $\Xx$ allows us to assume $f^{(0)}\leq f^\star$. Topicality then gives
\(f^{(0)}\leq f^{(1)}\leq\cdots\leq f^\star,\)
so the orbit converges pointwise.

The one-sided transforms inherit the modulus of continuity of the cost. Indeed, if
\(\delta(x,x')\eqdef\sup_{y\in\Yy}|c(x,y)-c(x',y)|\),
order reversal and additive anti-homogeneity give
\[
\left|g^{\bar c,\epsilon,\phi}(x)-g^{\bar c,\epsilon,\phi}(x')\right|
\leq \delta(x,x').
\]
Thus $(f^{(\ell)})_{\ell\geq1}$ is uniformly bounded and equicontinuous. Arzel\`a--Ascoli makes the orbit relatively compact in the uniform norm, and its monotonicity forces every uniformly convergent subsequence to have the same pointwise limit. Hence the full sequence converges uniformly.

Moreover, every selected one-sided transform is nonexpansive for the uniform norm. If $\norm{g-h}_\infty\leq r$, then $g-r\leq h\leq g+r$; order reversal and additive anti-homogeneity imply $\norm{g^{\bar c,\epsilon,\phi}-h^{\bar c,\epsilon,\phi}}_\infty\leq r$, and similarly for the other block. Therefore $\mathcal A_\phi$ is uniformly continuous, and the limit is a fixed point. The supersolution argument reverses all inequalities. Uniqueness modulo constants gives the last claim.
\end{proof}
\subsection{Sinkhorn Convergence: Sublinear Robust Rate}
\begin{prop}[Robust $O(1/\ell)$ dual rate for discrete Sinkhorn]\label{prop-sinkhorn-dual-rate}
Let $\a\in\simplex_n$ and $\b\in\simplex_m$ be positive histograms, and let $\C\in\RR_+^{n\times m}$ be finite.
Initialize $\gD^{(0)}=0$, perform complete Sinkhorn dual cycles, and denote the resulting potentials by $(\fD^{(\ell)},\gD^{(\ell)})$ for $\ell\geq1$. Let
\(\Delta^{(\ell)} \eqdef \Dd_\epsilon(\fD^\star,\gD^\star) - \Dd_\epsilon(\fD^{(\ell)},\gD^{(\ell)})\)
be the dual suboptimality gap for the entropic dual objective~\eqref{eq-dual-sinkhorn-objective}. Then
\(0\leq\Delta^{(\ell)}\leq \frac{2\norm{\C}_\infty^2}{\epsilon \ell}, \qquad \ell\geq1.\)
\end{prop}
\begin{proof}
Write $M\eqdef\norm{\C}_\infty$; the case $M=0$ is immediate because one complete cycle produces $\a\otimes\b$. We first prove the oscillation estimate used below. If $\fD$ is the soft transform of $\gD$, as in~\eqref{eq-discrete-soft-c-transforms}, set
\[
Z_i\eqdef\sum_j b_j\exp\!\left(\frac{g_j-C_{i,j}}{\epsilon}\right),
\qquad f_i=-\epsilon\log Z_i.
\]
Since $0\leq C_{i,j}\leq M$, one has $e^{-M/\epsilon}Z_{i'}\leq Z_i\leq e^{M/\epsilon}Z_{i'}$ for every $i,i'$. Therefore $\abs{f_i-f_{i'}}\leq M$ and $\norm{\fD}_V\leq M$. The same argument applies to the other soft transform. Every Sinkhorn potential is produced by such a transform, and the block optimality equations imply the same property for $(\fD^\star,\gD^\star)$.
\(\norm{\fD^{(\ell)}}_V,\ \norm{\gD^{(\ell)}}_V, \ \norm{\fD^\star}_V,\ \norm{\gD^\star}_V \leq M.\)

Associate with any potentials the nonnegative matrix
\[
\P(\fD,\gD)
\eqdef
(\a\otimes\b)
\odot
\exp\!\left(\frac{\fD\oplus\gD-\C}{\epsilon}\right).
\]
At the end of cycle $\ell$, the matrix $\P^{(\ell)}\eqdef\P(\fD^{(\ell)},\gD^{(\ell)})$ has column marginal $\b$ and hence total mass one. Denote its row marginal by $r^{(\ell)}\eqdef\P^{(\ell)}\ones_m$ and set $\delta_\ell\eqdef\norm{\a-r^{(\ell)}}_1$. The dual gradient at this point is $(\a-r^{(\ell)},0)$. Concavity of $\Dd_\epsilon$ therefore gives
\(\Delta^{(\ell)} \leq \dotp{\fD^\star-\fD^{(\ell)}}{\a-r^{(\ell)}}.\)
For every vector $h$ and every zero-sum vector $z$, subtracting $(\max h+\min h)/2$ from $h$ proves
\(\abs{\dotp{h}{z}} \leq \frac12\norm{h}_V\norm{z}_1.\)
The oscillation estimate and the triangle inequality for $\norm{\cdot}_V$ thus yield
\begin{equation}\label{eq-sinkhorn-gap-residual}
\Delta^{(\ell)}\leq M\delta_\ell.
\end{equation}

Its row update satisfies
\[
f_i^{(\ell+1)}-f_i^{(\ell)}
=
\epsilon\log\!\left(\frac{a_i}{r_i^{(\ell)}}\right).
\]
The matrices before and after this row update both have total mass one. Substitution in~\eqref{eq-dual-sinkhorn-objective} therefore gives the exact gain $\epsilon\KLD(\a|r^{(\ell)})$. The subsequent column update can only increase the dual objective, so Pinsker's inequality, Theorem~\ref{thm-pinsker}, and~\eqref{eq-sinkhorn-gap-residual} give
\[
\Delta^{(\ell)}-\Delta^{(\ell+1)}
\geq
\epsilon\KLD(\a|r^{(\ell)})
\geq
\frac{\epsilon}{2}\delta_\ell^2
\geq
\frac{\epsilon}{2M^2}(\Delta^{(\ell)})^2.
\]

Set $x=M/\epsilon$, and let $\widehat\P^{(1)}$ be the matrix after the first row update from $\gD^{(0)}=0$. Directly from the update, \(e^{-x} \leq \frac{\widehat P^{(1)}_{i,j}}{a_i b_j} \leq e^x.\)
Its column marginal $\widehat s$ consequently satisfies $e^{-x}\leq \widehat s_j/b_j\leq e^x$. The first column update gives $P^{(1)}_{i,j}=\widehat P^{(1)}_{i,j}b_j/\widehat s_j$, and hence
\(\frac{r_i^{(1)}}{a_i} = \sum_j\frac{\widehat P^{(1)}_{i,j}}{a_i}\frac{b_j}{\widehat s_j} \in[e^{-x},e^x],\)
because $(\widehat P^{(1)}_{i,j}/a_i)_j$ is a probability vector. Thus $\delta_1\leq\min\{2,e^x-1\}\leq2\min\{1,x\}$, where the last inequality uses $e^x-1\leq2x$ for $0\leq x\leq1$. Equation~\eqref{eq-sinkhorn-gap-residual} gives $\Delta^{(1)}\leq2M^2/\epsilon$.

Finally, whenever $\Delta^{(\ell+1)}>0$,
\[
\frac1{\Delta^{(\ell+1)}}-\frac1{\Delta^{(\ell)}}
=
\frac{\Delta^{(\ell)}-\Delta^{(\ell+1)}}{\Delta^{(\ell)}\Delta^{(\ell+1)}}
\geq
\frac{\epsilon}{2M^2}.
\]
Together with the bound on $\Delta^{(1)}$, telescoping proves $\Delta^{(\ell)}\leq2M^2/(\epsilon\ell)$. If a gap vanishes, all subsequent gaps vanish and the conclusion is immediate.
\end{proof}
\begin{cor}[Approximating unregularized OT by regularized dual costs]\label{cor-sinkhorn-dual-complexity}
Under the assumptions of Proposition~\ref{prop-sinkhorn-dual-rate}, suppose $nm>1$ and define the unregularized cost $L_0\eqdef\MKD_\C(\a,\b)$. For any $\delta>0$, choose $\epsilon=\delta/\log(nm)$. Then
\(\ell\geq \frac{2\norm{\C}_\infty^2\log(nm)}{\delta^2} \quad\Longrightarrow\quad \abs{L_0-\Dd_\epsilon(\fD^{(\ell)},\gD^{(\ell)})}\leq\delta.\)
\end{cor}
\begin{proof}
Let $L_\epsilon\eqdef\max_{\fD,\gD}\Dd_\epsilon(\fD,\gD)$ and let $\P^0$ solve the unregularized problem. The KL-normalized primal formulation~\eqref{eq-regularized-discr-rescaled} gives
\(0\leq L_\epsilon-L_0 \leq \epsilon\KLD(\P^0|\a\otimes\b) \leq \epsilon\log(nm).\)
The last inequality follows from $\KLD(\P^0|\a\otimes\b)=\HD(\a)+\HD(\b)-\HD(\P^0)\leq\log(nm)$.
Moreover, Proposition~\ref{prop-sinkhorn-dual-rate} gives $0\leq L_\epsilon-\Dd_\epsilon(\fD^{(\ell)},\gD^{(\ell)})\leq2\norm{\C}_\infty^2/(\epsilon\ell)$. Since the difference between the computed dual value and $L_0$ is the difference of these two nonnegative errors, its absolute value is bounded by their maximum. The prescribed $\epsilon$ and $\ell$ make both errors at most $\delta$.
\end{proof}
For dense $n\times n$ problems, one Sinkhorn cycle costs $O(n^2)$ operations. Corollary~\ref{cor-sinkhorn-dual-complexity} therefore yields the arithmetic complexity

\[
O\!\left(\frac{n^2\norm{\C}_\infty^2\log n}{\delta^2}\right)
\]
\subsection{Sinkhorn Convergence: Linear Hilbert Metric Rate}
\label{sec-sinkhorn-hilbert}
\paragraph{Projective contraction.}
\begin{defn}[Hilbert metric]\label{def-hilbert-metric}
On $\RR_{+,*}^n$, Hilbert's projective metric is
\eql{\label{eq-hilbert-metric}
\foralls (\uD,\uD') \in (\RR_{+,*}^n)^2, \quad
\Hilbert(\uD,\uD') \eqdef
\norm{\log(\uD)-\log(\uD')}_V.
}
Here, $\norm{\cdot}_V$ is the variation seminorm of Definition~\ref{def-variation-seminorm}, specialized to functions on the finite space $\{1,\ldots,n\}$: $\norm{z}_V=\max_i z_i-\min_i z_i$.
\end{defn}
\begin{prop}[Hilbert metric on the projective cone]
The function $\Hilbert$ defines a complete distance on the projective cone $\RR_{+,*}^n/\sim$, where $\uD \sim \uD'$ means that $\uD=s\uD'$ for some $s>0$.
\end{prop}
\begin{proof}
The map $\uD \mapsto \log \uD$ identifies $\RR_{+,*}^n/\sim$ with the quotient vector space $\RR^n/\Span(\ones_n)$, because multiplying $\uD$ by $s>0$ adds the constant vector $\log(s)\ones_n$. The variation seminorm $\norm{z}_V=\max_i z_i-\min_i z_i$ vanishes exactly on constant vectors, so it induces a norm on this quotient. Symmetry, the triangle inequality and separation therefore follow from the corresponding norm properties. Completeness follows because $\RR^n/\Span(\ones_n)$ is finite-dimensional and all finite-dimensional normed spaces are complete.
\end{proof}
\begin{thm}[Birkhoff contraction theorem]\label{thm-birkhoff}
Let $\K \in \RR_{+,*}^{n \times m}$. Then, for $(\vD,\vD') \in (\RR_{+,*}^m)^2$,
\[
\Hilbert(\K \vD,\K \vD') \leq \la(\K) \Hilbert(\vD,\vD'),
\qquad
\la(\K) \eqdef \frac{\sqrt{\eta(\K)}-1}{\sqrt{\eta(\K)}+1}<1,
\qquad
\eta(\K) \eqdef \umax{i,j,k,\ell}\frac{\K_{i,k}\K_{j,\ell}}{\K_{j,k}\K_{i,\ell}}.
\]
\end{thm}
\begin{proof}
Work in logarithmic coordinates and define
\(F(z)\eqdef\log(\K e^z), \qquad \mathsf P(z)_{i,k} \eqdef \frac{\K_{i,k}e^{z_k}}{(\K e^z)_i}.\)
Every row of $\mathsf P(z)$ is a probability vector and $\d F(z)[h]=\mathsf P(z)h$. We first bound its contraction in variation seminorm. If $p$ and $q$ are two probability rows, then
\(\abs{\dotp{p-q}{h}} \leq \frac12\norm{p-q}_1\norm{h}_V.\)
Indeed, subtracting $\min_k h_k$ and dividing by $\norm{h}_V$ reduces this inequality to the variational characterization of total variation by functions taking values in $[0,1]$.
\begin{equation}\label{eq-birkhoff-proof-dobrushin}
\norm{\mathsf P(z)h}_V
\leq
\delta(\mathsf P(z))\norm{h}_V,
\qquad
\delta(\mathsf P)
\eqdef
\frac12\max_{i,j}\sum_k\abs{\mathsf P_{i,k}-\mathsf P_{j,k}}.
\end{equation}

Fix rows $p,q$ of $\mathsf P(z)$, with indices $i,j$, and put $r_k=p_k/q_k$. Then $\sum_kq_kr_k=1$, and the definition of $\eta(\K)$ gives
\(\frac{\max_k r_k}{\min_k r_k} = \max_{k,\ell} \frac{\K_{i,k}\K_{j,\ell}}{\K_{j,k}\K_{i,\ell}} \leq\eta(\K).\)
Write $a=\min_k r_k\leq1\leq b=\max_k r_k$. If $a=b=1$, then $p=q$ and the desired bound is immediate. Otherwise, $(r-1)_+\leq (b-1)(r-a)/(b-a)$ on $[a,b]$, so
\(\frac12\norm{p-q}_1 = \sum_kq_k(r_k-1)_+ \leq \frac{(1-a)(b-1)}{b-a}.\)
Set $a=e^{-u}$ and $b=e^v$. The last quotient equals
\[
\frac{2\sinh(u/2)\sinh(v/2)}{\sinh((u+v)/2)}
\leq
\tanh\!\left(\frac{u+v}{4}\right)
\leq
\tanh\!\left(\frac{\log\eta(\K)}{4}\right)
=
\la(\K),
\]
where the first inequality follows from $\sinh(u/2)\sinh(v/2)\leq\sinh^2((u+v)/4)$. Thus~\eqref{eq-birkhoff-proof-dobrushin} holds with $\delta(\mathsf P(z))\leq\la(\K)$, uniformly in $z$.

For $z_t=(1-t)z'+tz$, the fundamental theorem of calculus and the triangle inequality for $\norm{\cdot}_V$ now give
\(\norm{F(z)-F(z')}_V \leq \int_0^1\norm{\mathsf P(z_t)(z-z')}_V\d t \leq \la(\K)\norm{z-z'}_V.\)
Taking $z=\log\vD$ and $z'=\log\vD'$ proves the result because $F(z)=\log(\K\vD)$.
\end{proof}
\(\Hilbert(\K^\ell u^{(0)},u^\star) \leq \la(\K)^\ell\Hilbert(u^{(0)},u^\star).\)
Thus Hilbert contraction supplies both the Perron--Frobenius eigenvector and the global linear convergence of power iteration. The contraction can be visualized for every initialization at once on the three-state probability simplex.

\(\K^{\ell+1}\simplex_3 = \K^\ell(\K\simplex_3) \subseteq \K^\ell\simplex_3.\)
\[
J_3\eqdef\frac{1}{3}\ones_3\ones_3^\top,
\qquad
A\eqdef
\begin{pmatrix}
0&-1&1\\
1&0&-1\\
-1&1&0
\end{pmatrix},
\qquad
\K_{\rho,\delta}\eqdef \rho I_3+(1-\rho)J_3+\delta A.
\]
\[
\K_1=\K_{.90,0},
\qquad
\K_2=\K_{.90,.014},
\qquad
\K_3=
\begin{pmatrix}
.950&.018&.032\\
.042&.880&.078\\
.008&.102&.890
\end{pmatrix}.
\]
\paragraph{Sinkhorn contraction.}
\begin{thm}[Projective linear convergence of Sinkhorn]\label{thm-sinkhorn-hilbert-linear}
Starting from $\vD^{(0)}>0$, use the scaling iterates of~\eqref{eq-sinkhorn} and the associated primal half-step iterates already defined in~\eqref{eq-kl-sinkh-proj}. Fix any optimal scaling pair $(\uD^\star,\vD^\star)$.
Writing $\la=\la(\K)$, one has
\eql{\label{eq-convlin-sinkh}
\Hilbert(\vD^{(\ell)},\vD^\star)
\leq
\la^{2\ell}\Hilbert(\vD^{(0)},\vD^\star),
\qquad
\Hilbert(\uD^{(\ell+1)},\uD^\star)
\leq
\la^{2\ell+1}\Hilbert(\vD^{(0)},\vD^\star).
}
Thus the scaling rays converge linearly. After fixing any continuous gauge, their representatives converge, and $\P^{(\ell)}\to\P^\star$ entrywise.
\eql{\label{eq-convsinkh-control}
\Hilbert(\uD^{(\ell)},\uD^\star)
\leq
\frac{\Hilbert(\P^{(\ell)}\ones_m,\a)}{1-\la^2}
\qandq
\Hilbert(\vD^{(\ell)},\vD^\star)
\leq
\frac{\Hilbert((\P^{(\ell+1/2)})^\top\ones_n,\b)}{1-\la^2}.
}
\eqllead{Lastly,}{eq-convlin-sinkh-prim}{
\norm{\log(\P^{(\ell)})-\log(\P^\star)}_\infty
\leq
\Hilbert(\uD^{(\ell)},\uD^\star)
+
\Hilbert(\vD^{(\ell)},\vD^\star),
}
where $\P^\star$ is the unique solution of~\eqref{eq-regularized-discr}.
\end{thm}
\begin{proof}
Entrywise multiplication by a fixed positive vector and entrywise inversion are isometries of Hilbert's metric. Therefore the row map $R(\vD)=\a/(\K\vD)$ and the column map $C(\uD)=\b/(\K^\top\uD)$ are both $\la$-Lipschitz by Theorem~\ref{thm-birkhoff}. Their full-cycle compositions $C\circ R$ and $R\circ C$ are $\la^2$-contractions. Since $\vD^\star=C(R(\vD^\star))$ and $\uD^\star=R(C(\uD^\star))$, iteration gives~\eqref{eq-convlin-sinkh}.

For any contraction $F$ with factor $q<1$ and fixed point $x^\star$, the triangle inequality gives \(d(x,x^\star) \leq d(x,Fx)+d(Fx,Fx^\star) \leq d(x,Fx)+q\,d(x,x^\star),\)
hence $d(x,x^\star)\leq d(x,Fx)/(1-q)$. Apply this with $q=\la^2$ to the two full-cycle maps. The scaling identities
\(
\Hilbert(\uD^{(\ell+1)},\uD^{(\ell)})
=\Hilbert(\a,\P^{(\ell)}\ones_m)
\)
and
\(
\Hilbert(\vD^{(\ell+1)},\vD^{(\ell)})
=\Hilbert(\b,(\P^{(\ell+1/2)})^\top\ones_n)
\)
then prove~\eqref{eq-convsinkh-control}.

Finally, write $\xi_i=\log(u_i^{(\ell)}/u_i^\star)$ and $\zeta_j=\log(v_j^{(\ell)}/v_j^\star)$. Then $\log(\P_{i,j}^{(\ell)}/\P_{i,j}^\star)=\xi_i+\zeta_j$ and $\operatorname{osc}(\xi\oplus\zeta)=\operatorname{osc}(\xi)+\operatorname{osc}(\zeta)$.
Both plans have total mass one, so $\sum_{i,j}\P_{i,j}^\star e^{\xi_i+\zeta_j}=1$. This proves~\eqref{eq-convlin-sinkh-prim}.
\end{proof}
\paragraph{Nonlinear Sinkhorn images of the simplex.}
\[
F_{\uD}(\uD)
\eqdef
R(C(\uD))
=
\a\oslash\left[\K\left(\b\oslash(\K^\top\uD)\right)\right].
\]
\(\widehat F_{\uD}(p) \eqdef \frac{F_{\uD}(p)}{\dotp{F_{\uD}(p)}{\ones_3}}, \qquad p\in\simplex_3.\)
Thus $\widehat\uD^{(\ell)}\eqdef\uD^{(\ell)}/\dotp{\uD^{(\ell)}}{\ones_3}$ follows $\widehat\uD^{(\ell+1)}=\widehat F_{\uD}(\widehat\uD^{(\ell)})$ when $\ell$ indexes complete Sinkhorn cycles. Since $\widehat F_{\uD}$ maps $\simplex_3$ into itself, its set-valued iterates are nested:

\[
\widehat F_{\uD}^{\,\ell+1}(\simplex_3)
=
\widehat F_{\uD}^{\,\ell}\!\left(\widehat F_{\uD}(\simplex_3)\right)
\subseteq
\widehat F_{\uD}^{\,\ell}(\simplex_3).
\]
\[
\operatorname{diam}_{\Hilbert}\!\left(\widehat F_{\uD}^{\,\ell}(\simplex_3)\right)
\leq
\la(\K)^{2(\ell-1)}
\operatorname{diam}_{\Hilbert}\!\left(\widehat F_{\uD}(\simplex_3)\right),
\qquad \ell\geq1.
\]
\paragraph{Dual-potential form of the contraction.}
\[
\norm{\fD^{(\ell)}-\fD^\star}_V
=
\epsilon\Hilbert(\uD^{(\ell)},\uD^\star),
\qquad
\norm{\gD^{(\ell)}-\gD^\star}_V
=
\epsilon\Hilbert(\vD^{(\ell)},\vD^\star).
\]
Thus a complete cycle contracts both potential classes at the rate $\la(\K)^2$. Moreover, the temperature factor must be retained when passing from potentials to densities:

\eq{
\norm{ \log \frac{\d\pi^{(\ell)}}{\d\pi^\star} }_\infty
=
\frac1\epsilon
\norm{ (\fD^{(\ell)}-\fD^\star) \oplus (\gD^{(\ell)}-\gD^\star) }_\infty
\leq
\frac{\norm{\fD^{(\ell)}-\fD^\star}_V + \norm{\gD^{(\ell)}-\gD^\star}_V}{\epsilon}.
}
\(\la=\frac{\sqrt{\eta}-1}{\sqrt{\eta}+1} \leq \tanh(R/(2\epsilon))<1.\)
\(\norm{\fD^{(\ell)}-\fD^\star}_V \leq \frac{\epsilon}{1-\la^2} \norm{\log((\P^{(\ell)}\ones_m)\oslash\a)}_V,\)
\(\norm{\gD^{(\ell)}-\gD^\star}_V \leq \frac{\epsilon}{1-\la^2} \norm{\log(((\P^{(\ell+1/2)})^\top\ones_n)\oslash\b)}_V.\)
\subsection{Local Convergence Analysis and Acceleration}
\label{sec-sinkhorn-local-acceleration}
\paragraph{Conditional operator and maximal correlation.}
\begin{defn}[Conditional coupling operator]\label{def-sinkhorn-conditional-operator}
For a coupling $\pi\in\Couplings(\al,\be)$, let
\(
\pi(\d x,\d y)=\al(\d x)\pi_x(\d y)
=\be(\d y)\pi^y(\d x)
\)
be its two disintegrations. Define
\(
T_\pi:L^2(\be)\to L^2(\al)
\)
and
\(
T_\pi^*:L^2(\al)\to L^2(\be)
\)
by
\eql{\label{eq-sinkhorn-conditional-operator}
(T_\pi k)(x)\eqdef\int_\Yy k(y)\d\pi_x(y),
\qquad
(T_\pi^*h)(y)\eqdef\int_\Xx h(x)\d\pi^y(x).
}
The centered maximal correlation of $\pi$ is
\eql{\label{eq-sinkhorn-maximal-correlation}
\sigma(\pi)
\eqdef
\sup_{\substack{k\in L_0^2(\be)\\\norm{k}_{L^2(\be)}=1}}
\norm{T_\pi k}_{L^2(\al)}
=
\sup_{\substack{h\in L_0^2(\al),\,k\in L_0^2(\be)\\
\norm{h}_{L^2(\al)}=\norm{k}_{L^2(\be)}=1}}
\int h(x)k(y)\d\pi(x,y).
}
For the optimal entropic coupling $\pi_\epsilon$ of~\eqref{eq-entropic-generic}, set
\(
T_\epsilon\eqdef T_{\pi_\epsilon}
\)
and
\(
\sigma_\epsilon\eqdef\sigma(\pi_\epsilon).
\)
\end{defn}
\eql{\label{eq-sinkhorn-conditional-operator-discrete}
T_\epsilon=\diag(\a)^{-1}\P_\epsilon,
\qquad
T_\epsilon^*=\diag(\b)^{-1}\P_\epsilon^\top,
}
where adjointness is understood in the $\a$- and $\b$-weighted Euclidean inner products. Thus $\sigma_\epsilon$ is the largest nonconstant singular value of the normalized optimal coupling.

\paragraph{Dual Hessian and exact local rate.}
\begin{prop}[Dual Hessian and local Sinkhorn rate]\label{prop-sinkhorn-local-rate}
Assume the setting of Proposition~\ref{prop-continuous-entropic-duality}, and that differentiation under the Gibbs integrals and the stated $L^2$ Taylor expansions are valid. These requirements are automatic in finite spaces and form the local regularity hypothesis in continuous settings. For bounded directions $(h,k)$ and $(h',k')$,
\eql{\label{eq-sinkhorn-dual-hessian-form}
-\epsilon D^2\Dd_\epsilon(f_\epsilon,g_\epsilon)
[(h,k),(h',k')]
=
\int_{\Xx\times\Yy}
(h(x)+k(y))(h'(x)+k'(y))\d\pi_\epsilon(x,y).
}
Hence the Hessian form extends continuously to
\(L^2(\al)\oplus L^2(\be)\), where the positive dual curvature operator is
\eql{\label{eq-sinkhorn-dual-hessian-operator}
-\nabla^2\Dd_\epsilon(f_\epsilon,g_\epsilon)
=
\frac1\epsilon
\begin{pmatrix}
\Id&T_\epsilon\\
T_\epsilon^*&\Id
\end{pmatrix}.
}
After quotienting by the gauge direction $(1,-1)$ and using the product-$L^2$ quotient norm, its spectrum is contained in $[(1-\sigma_\epsilon)/\epsilon,2/\epsilon]$.
Both endpoints belong to the spectrum, although the lower one need not be an eigenvalue.
In particular, when $\sigma_\epsilon<1$, the spectral condition number of the full dual is \(2/(1-\sigma_\epsilon)\).
The derivatives of the two soft-transform updates~\eqref{eq-soft-c-cont-f}--\eqref{eq-soft-c-cont-g} at the optimum are
\eql{\label{eq-sinkhorn-soft-transform-jacobians}
D[g\mapsto g^{\bar c,\epsilon}](g_\epsilon)=-T_\epsilon,
\qquad
D[f\mapsto f^{c,\epsilon}](f_\epsilon)=-T_\epsilon^*.
}
\(g\mapsto(g^{\bar c,\epsilon})^{c,\epsilon}\)
is $T_\epsilon^*T_\epsilon$. If the gauges are fixed so that
\(e^{(\ell)}\eqdef g^{(\ell)}-g_\epsilon\in L_0^2(\be)\), then
\eql{\label{eq-sinkhorn-local-error}
e^{(\ell+1)}
=
T_\epsilon^*T_\epsilon e^{(\ell)}
+
o\!\left(\norm{e^{(\ell)}}_{L^2(\be)}\right),
\qquad
\limsup_{\ell\to\infty}
\frac{\norm{e^{(\ell+1)}}_{L^2(\be)}}
{\norm{e^{(\ell)}}_{L^2(\be)}}
\leq\sigma_\epsilon^2.
}
If $\sigma_\epsilon^2$ is an isolated eigenvalue and the asymptotic error has a nonzero projection onto its eigenspace, the ratio in~\eqref{eq-sinkhorn-local-error} converges to $\sigma_\epsilon^2$.
\end{prop}
\begin{proof}
Twice differentiating the exponential term in~\eqref{eq-dual-sinkhorn-objective} gives~\eqref{eq-sinkhorn-dual-hessian-form}; the density law~\eqref{eq-continuous-entropic-density-law} identifies its reference measure with $\pi_\epsilon$. Its two diagonal terms are the $L^2$ inner products of the marginals, and its cross term is represented by $T_\epsilon$, which proves~\eqref{eq-sinkhorn-dual-hessian-operator}.

Each quotient class has a unique representative orthogonal to the gauge direction, hence of the form
\(
(h_0+m,k_0+m)
\)
with $h_0\in L_0^2(\al)$ and $k_0\in L_0^2(\be)$. On the centered component, Cauchy--Schwarz and~\eqref{eq-sinkhorn-maximal-correlation} give
\[
(1-\sigma_\epsilon)
\bigl(\norm{h_0}_{L^2(\al)}^2+\norm{k_0}_{L^2(\be)}^2\bigr)
\leq
\norm{h_0}_{L^2(\al)}^2+\norm{k_0}_{L^2(\be)}^2
+2\langle h_0,T_\epsilon k_0\rangle_{L^2(\al)}
\leq
(1+\sigma_\epsilon)
\bigl(\norm{h_0}_{L^2(\al)}^2+\norm{k_0}_{L^2(\be)}^2\bigr).
\]
The common-constant direction $(1,1)$ has eigenvalue $2/\epsilon$, whereas $(1,-1)$ is the gauge kernel. Approximate maximizers in~\eqref{eq-sinkhorn-maximal-correlation} show that the lower spectral endpoint is $(1-\sigma_\epsilon)/\epsilon$, while the constant direction gives the upper endpoint $2/\epsilon$. This proves the condition-number formula without requiring a singular-vector basis.

Differentiating either log-partition formula in Definition~\ref{def-continuous-soft-c-transform} produces minus the corresponding conditional expectation under $\pi_\epsilon$, proving~\eqref{eq-sinkhorn-soft-transform-jacobians}. The chain rule and a Taylor expansion then give~\eqref{eq-sinkhorn-local-error}. Since
\(\norm{T_\epsilon^*T_\epsilon}_{L_0^2(\be)\to L_0^2(\be)}
=\sigma_\epsilon^2\),
the rate follows.
\end{proof}
\eql{\label{eq-sinkhorn-conditional-variance-gap}
1-\sigma_\epsilon^2
=
\inf_{\substack{k\in L_0^2(\be)\\\norm{k}_{L^2(\be)}=1}}
\int_\Xx
\operatorname{Var}_{\pi_{\epsilon,x}}\!\bigl(k(Y)\bigr)
\d\al(x).
}
Thus Sinkhorn slows down when the optimal conditional laws $Y|X=x$ become nearly deterministic. This is exactly what happens as $\epsilon\to0$ in a smooth Monge regime.

\begin{prop}[The global Hilbert factor controls the local rate]\label{prop-sinkhorn-hilbert-controls-local}
In the discrete positive-kernel setting of Theorem~\ref{thm-sinkhorn-hilbert-linear}, let $\la(\K)$ be the Birkhoff contraction factor defined in Theorem~\ref{thm-birkhoff}. Then
\eql{\label{eq-sinkhorn-hilbert-controls-local}
\sigma_\epsilon\leq\la(\K)<1,
\qquad
\sigma_\epsilon^2\leq\la(\K)^2.
}
\end{prop}
\begin{proof}
Let $R(\vD)=\a\oslash(\K\vD)$ and $C(\uD)=\b\oslash(\K^\top\uD)$ be the row and column scaling maps used in the proof of Theorem~\ref{thm-sinkhorn-hilbert-linear}. Entrywise inversion and multiplication by a fixed positive vector are isometries of Hilbert's metric. Theorem~\ref{thm-birkhoff}, applied successively to $\K$ and $\K^\top$, therefore gives the pairwise estimate
\(\Hilbert\bigl(C(R(\vD)),C(R(\widetilde\vD))\bigr) \leq \la(\K)^2\Hilbert(\vD,\widetilde\vD),\)
where $\la(\K^\top)=\la(\K)$ follows directly from its cross-ratio definition. Under the logarithmic change of variables $g=\epsilon\log\vD$, Hilbert's metric becomes $\epsilon^{-1}\norm{\cdot}_V$ by~\eqref{eq-hilbert-metric}. Hence the complete dual update
\(
\mathcal S(g)\eqdef(g^{\bar c,\epsilon})^{c,\epsilon}
\)
is $\la(\K)^2$-Lipschitz for the variation norm on the quotient by constants: \(\norm{\mathcal S(g)-\mathcal S(\widetilde g)}_V \leq \la(\K)^2\norm{g-\widetilde g}_V.\)
Differentiate this inequality at the fixed point $g_\epsilon$. Proposition~\ref{prop-sinkhorn-local-rate} gives
\(D\mathcal S(g_\epsilon)=T_\epsilon^*T_\epsilon\), and therefore
\(\norm{T_\epsilon^*T_\epsilon h}_V \leq \la(\K)^2\norm{h}_V \qquad \text{for every }h\in L_0^2(\be).\)
In the present finite-dimensional setting, $T_\epsilon^*T_\epsilon$ is positive and self-adjoint for the $\be$-weighted inner product, and its largest eigenvalue on $L_0^2(\be)$ is $\sigma_\epsilon^2$. If $\sigma_\epsilon>0$, take a corresponding eigenvector $h_\star$. It is nonconstant, so $\norm{h_\star}_V>0$, and the preceding bound gives
\(\sigma_\epsilon^2\norm{h_\star}_V = \norm{T_\epsilon^*T_\epsilon h_\star}_V \leq \la(\K)^2\norm{h_\star}_V.\)
The conclusion is immediate; the case $\sigma_\epsilon=0$ is trivial.
\end{proof}
\paragraph{Blockwise over-relaxation.}
\eql{\label{eq-sinkhorn-block-overrelaxation}
f^{(\ell+1)}=(1-\omega)f^{(\ell)}+\omega(g^{(\ell)})^{\bar c,\epsilon},
\qquad
g^{(\ell+1)}=(1-\omega)g^{(\ell)}+\omega(f^{(\ell+1)})^{c,\epsilon}.
}
\begin{prop}[Optimal local block relaxation]\label{prop-sinkhorn-optimal-overrelaxation}
Assume the hypotheses of Proposition~\ref{prop-sinkhorn-local-rate} and $\sigma_\epsilon<1$. For every $0<\omega<2$, the iteration~\eqref{eq-sinkhorn-block-overrelaxation} is locally linearly convergent modulo the additive gauge. Its optimal relaxation parameter and corresponding asymptotic factor are
\eql{\label{eq-sinkhorn-optimal-overrelaxation}
\omega_\star
=
\frac{2}{1+\sqrt{1-\sigma_\epsilon^2}},
\qquad
r_\star
=
\omega_\star-1
=
\frac{1-\sqrt{1-\sigma_\epsilon^2}}
{1+\sqrt{1-\sigma_\epsilon^2}}.
}
For $\omega=1$, the factor is $\sigma_\epsilon^2$, as in Proposition~\ref{prop-sinkhorn-local-rate}.
\end{prop}
\begin{proof}
Use~\eqref{eq-sinkhorn-soft-transform-jacobians}. The polar decomposition of $T_\epsilon$ and the spectral theorem for $(T_\epsilon^*T_\epsilon)^{1/2}$ reduce the centered linearized iteration to scalar spectral values $s\in[0,\sigma_\epsilon]$; in finite dimensions these are the usual singular values. On the corresponding two-dimensional mode, one relaxed cycle has matrix
\[
\mathsf M_\omega(s)
=
\begin{pmatrix}
1-\omega&-\omega s\\
-\omega s(1-\omega)&1-\omega+\omega^2s^2
\end{pmatrix},
\]
whose characteristic polynomial is $z^2-(2(1-\omega)+\omega^2s^2)z+(1-\omega)^2$.
The supremum of the root moduli is reached, or approached, as $s\to\sigma_\epsilon$. The two roots of this limiting mode have equal modulus precisely when
\(\omega^2\sigma_\epsilon^2=4(\omega-1)\).
The smaller solution is $\omega_\star$ in~\eqref{eq-sinkhorn-optimal-overrelaxation}; at this value both root moduli equal $\omega_\star-1$. Below this threshold the dominant real root decreases with $\omega$, while above it the conjugate-root modulus $\omega-1$ increases. This proves optimality and local convergence for $0<\omega<2$.
The nongauge constant mode has factor $(1-\omega)^2$ and never dominates these centered modes.
\end{proof}
Set $\delta_\epsilon=1-\sigma_\epsilon^2$. When $\delta_\epsilon$ is small, the exact expression in~\eqref{eq-sinkhorn-optimal-overrelaxation} gives the complete-cycle factor

\(r_\star = \frac{1-\sqrt{\delta_\epsilon}}{1+\sqrt{\delta_\epsilon}} = 1-2\sqrt{\delta_\epsilon} +2\delta_\epsilon +O(\delta_\epsilon^{3/2}),\)
\paragraph{Entropic semi-dual and variable projection.}
The hard semi-dual \(\Ee_0\) in~\eqref{eq-semi-dual} eliminates one potential from the full dual \(\Dd_0\). At positive temperature, the full dual is \(\Dd_\epsilon\) from~\eqref{eq-dual-sinkhorn-objective}, and the soft \(\bar c\)-transform \(g^{\bar c,\epsilon}\) is defined by~\eqref{eq-soft-c-cont-g}.

\eql{\label{eq-entropic-semidual-local}
\Ee_\epsilon(g)
\eqdef
\Dd_\epsilon(g^{\bar c,\epsilon},g)
=
\int_\Xx g^{\bar c,\epsilon}\d\al
+
\int_\Yy g\d\be.
}
For a potential $g$, let $\pi_g$ be the Gibbs measure obtained from the density law~\eqref{eq-continuous-entropic-density-law} by replacing $(f_\epsilon,g_\epsilon)$ with $(g^{\bar c,\epsilon},g)$. Write \( \pi_g(\d x,\d y)=\al(\d x)\pi_{g,x}(\d y) \) and denote its second marginal by $\be_g$. The first marginal is exactly $\al$ by construction.

\(\min_{U,V}\frac12\norm{U\sigma(VX)-Y}_{\rm F}^2,\)
\begin{prop}[Semi-dual curvature and conditioning]\label{prop-sinkhorn-semidual-curvature}
For bounded directions $k,k'$,
\eql{\label{eq-entropic-semidual-derivatives}
D\Ee_\epsilon(g)[k]
=
\int_\Yy k\d(\be-\be_g),
\qquad
D^2\Ee_\epsilon(g)[k,k']
=
-\frac1\epsilon
\int_\Xx
\operatorname{Cov}_{\pi_{g,x}}\!\bigl(k(Y),k'(Y)\bigr)
\d\al(x).
}
At $g=g_\epsilon$, its positive curvature on $L_0^2(\be)$ is
\eql{\label{eq-entropic-semidual-hessian}
-\nabla^2\Ee_\epsilon(g_\epsilon)
=
\frac1\epsilon(\Id-T_\epsilon^*T_\epsilon),
\qquad
\frac{1-\sigma_\epsilon^2}{\epsilon}\Id
\preceq
-\nabla^2\Ee_\epsilon(g_\epsilon)
\preceq
\frac1\epsilon\Id.
}
Hence its spectral condition number is at most
\(
1/(1-\sigma_\epsilon^2)
\).
\end{prop}
\begin{proof}
The envelope theorem gives the first derivative in~\eqref{eq-entropic-semidual-derivatives}. Differentiating the normalized conditional Gibbs law gives the covariance formula. At the optimum, total covariance yields
\(\int_\Xx \operatorname{Cov}_{\pi_{\epsilon,x}}(k(Y),k'(Y)) \d\al(x) = \int_\Yy kk'\d\be - \int_\Xx(T_\epsilon k)(T_\epsilon k')\d\al,\)
which is~\eqref{eq-entropic-semidual-hessian}. Equivalently, this operator is the Schur complement of the first identity block in~\eqref{eq-sinkhorn-dual-hessian-operator}. Its spectrum on centered functions lies in
\(
[(1-\sigma_\epsilon^2)/\epsilon,1/\epsilon]
\).
\end{proof}
\subsection{Entropic Optimal Transport between Gaussians}
\label{sec-gaussian-sinkhorn}
\begin{prop}[Quadratic closure of Sinkhorn iterates]\label{prop-gaussian-sinkhorn-closure}
Let $\be=\Gaussian(\mean_\be,\cov_\be)$ on $\RR^d$ and take $c(x,y)=\norm{x-y}^2$. If $g(y)$ is a quadratic polynomial such that the Gaussian integral below is finite, then the soft transform
\[
f(x)
=
-\epsilon\log
\int
\exp\!\left(\frac{g(y)-\norm{x-y}^2}{\epsilon}\right)
\d\be(y)
\]
is a quadratic polynomial in $x$.
\end{prop}
\begin{proof}
The exponent is the sum of a quadratic polynomial in $y$ and the logarithm of the Gaussian density of $\be$. Completing the square in $y$ evaluates the integral as a positive constant times the exponential of a quadratic polynomial in $x$. Taking $-\epsilon\log$ therefore gives a quadratic polynomial.
\end{proof}
\begin{prop}[Balanced entropic OT between Gaussians]\label{prop-gaussian-sinkhorn-closed-form}
Let $\al=\Gaussian(\mean_\al,\cov_\al)$ and $\be=\Gaussian(\mean_\be,\cov_\be)$ have positive-definite covariances, and let $\cov_\al^{1/2}\cov_\be^{1/2}=U\diag(\sigma_i)V^\top$ be a singular-value decomposition. For the balanced entropic problem~\eqref{eq-entropic-generic} with $c(x,y)=\norm{x-y}^2$,
the optimizer is Gaussian with cross-covariance
\(K_\epsilon = \cov_\al^{1/2} U\diag(s_i)V^\top \cov_\be^{1/2}, \qquad s_i = \frac{\sqrt{\epsilon^2+16\sigma_i^2}-\epsilon}{4\sigma_i}.\)
The optimal value is
\[
\norm{\mean_\al-\mean_\be}^2
+
\tr(\cov_\al)+\tr(\cov_\be)
+
\sum_i
\left(
-2\sigma_i s_i
-\frac{\epsilon}{2}\log(1-s_i^2)
\right).
\]
As $\epsilon\downarrow0$, $s_i\to1$ and the full covariance contribution, including the two trace terms, converges to $\Bb(\cov_\al,\cov_\be)^2$.
\end{prop}
\begin{proof}
Let $(X,Y)$ be any coupling with finite second moments and cross-covariance
\(K=\EE\big[(X-\mean_\al)(Y-\mean_\be)^\top\big]\).
Replacing $(X,Y)$ by the Gaussian vector with the same mean and covariance leaves the quadratic cost unchanged. Since the marginals are fixed, \(\KL(\pi|\al\otimes\be) = -h(X,Y)+h(\al)+h(\be),\)
where $h$ denotes differential entropy when it is finite. Among laws with a fixed covariance, the Gaussian maximizes entropy; if the entropy is not finite, the relative entropy is already $+\infty$. Thus the Gaussian replacement cannot increase the objective, and it is enough to optimize over Gaussian couplings.

Any such coupling has covariance
\[
\begin{pmatrix}
\cov_\al & K\\
K^\top & \cov_\be
\end{pmatrix}.
\]
Write $K=\cov_\al^{1/2}S\cov_\be^{1/2}$. The block covariance constraint is equivalent to the singular values of $S$ being at most one, and finite entropy forces them to be strictly smaller than one. The cost depends on $K$ through
\(\norm{\mean_\al-\mean_\be}^2+\tr(\cov_\al)+\tr(\cov_\be)-2\tr(K),\)
while $\KL(\pi|\al\otimes\be)=-\frac12\log\det(\Id-SS^\top)$.
By von Neumann's trace inequality, the minimizer aligns $S$ with the singular vectors of $\cov_\al^{1/2}\cov_\be^{1/2}$, so $S=U\diag(s_i)V^\top$. The problem separates into scalar minimizations
\(\min_{0\leq s<1} -2\sigma_i s-\frac{\epsilon}{2}\log(1-s^2).\)
The first-order condition is $2\sigma_i=\epsilon s/(1-s^2)$, whose positive solution is the displayed $s_i$. Substitution gives the value formula. Since $s_i\to1$ and $\epsilon\log(1-s_i^2)\to0$ as $\epsilon\downarrow0$, the spectral sum converges to $-2\sum_i\sigma_i$. The identity $\sum_i\sigma_i=\tr((\cov_\al^{1/2}\cov_\be\cov_\al^{1/2})^{1/2})$ then gives the Bures--Wasserstein covariance contribution.
\end{proof}
\begin{cor}[Gaussian Sinkhorn divergence and smoothed Bures term]\label{cor-gaussian-sinkhorn-divergence}
Let $\al=\Gaussian(\mean_\al,\cov_\al)$ and $\be=\Gaussian(\mean_\be,\cov_\be)$ have positive-definite covariances. For $r>0$, define
\[
\tau_\epsilon(r)
\eqdef
\frac{\sqrt{\epsilon^2+16r^2}-\epsilon}{4r},
\qquad
\psi_\epsilon(r)
\eqdef
-2r\,\tau_\epsilon(r)
-\frac{\epsilon}{2}\log\bigl(1-\tau_\epsilon(r)^2\bigr).
\]
If $\sigma_i(\Sigma,\Lambda)$ denotes the singular values of $\Sigma^{1/2}\Lambda^{1/2}$ and $\lambda_i(\Sigma)$ the eigenvalues of $\Sigma$, then the debiased Sinkhorn divergence~\eqref{eq-sinkhorn-divergence} is
\(\bar\MK_{\norm{\cdot-\cdot}^2}^{\epsilon}(\al,\be) = \norm{\mean_\al-\mean_\be}^2 + \Bb_\epsilon(\cov_\al,\cov_\be)^2,\)
where the Gaussian covariance contribution is the debiased smoothed Bures term
\[
\Bb_\epsilon(\Sigma,\Lambda)^2
\eqdef
\sum_i \psi_\epsilon\bigl(\sigma_i(\Sigma,\Lambda)\bigr)
-\frac12\sum_i\psi_\epsilon\bigl(\lambda_i(\Sigma)\bigr)
-\frac12\sum_i\psi_\epsilon\bigl(\lambda_i(\Lambda)\bigr).
\]
Moreover $\Bb_\epsilon(\Sigma,\Lambda)^2\to\Bb(\Sigma,\Lambda)^2$ as $\epsilon\downarrow0$, where $\Bb$ is the Bures--Wasserstein metric of Proposition~\ref{prop-gaussian-w2-bures}.
\end{cor}
\begin{proof}
Proposition~\ref{prop-gaussian-sinkhorn-closed-form} writes the raw entropic value as the squared mean displacement plus trace terms and a spectral sum. With the notation above, the spectral part is exactly $\sum_i\psi_\epsilon(\sigma_i(\cov_\al,\cov_\be))$. Applying the same formula to the self-costs $(\al,\al)$ and $(\be,\be)$ replaces these singular values by the eigenvalues of $\cov_\al$ and $\cov_\be$. In the polarization formula~\eqref{eq-sinkhorn-divergence}, the trace terms cancel because
$\tr\cov_\al+\tr\cov_\be-\frac12(2\tr\cov_\al)-\frac12(2\tr\cov_\be)=0$, while the polarization of the squared mean terms leaves exactly $\norm{\mean_\al-\mean_\be}^2$. Since $\tau_\epsilon(r)\to1$ and $\epsilon\log(1-\tau_\epsilon(r)^2)\to0$, one has $\psi_\epsilon(r)\to-2r$. The limit is therefore
\(\tr\Sigma+\tr\Lambda -2\sum_i\sigma_i(\Sigma,\Lambda) = \Bb(\Sigma,\Lambda)^2,\)
which is the Bures formula~\eqref{eq-bure-defn}.
\end{proof}
\begin{prop}[One-dimensional Gaussian Sinkhorn rate]\label{prop-gaussian-sinkhorn-1d-rate}
Consider $\al=\be=\Gaussian(0,1)$ on $\RR$ with $c(x,y)=(x-y)^2$. Initialize $g^{(0)}=0$ and, after $\ell$ complete Sinkhorn cycles, write
\(g^{(\ell)}(y)=q^{(\ell)}y^2+\mathrm{cst}.\)
One soft transform maps a quadratic coefficient $q$ to
\(\mathsf Q_\epsilon(q) = 1-\frac{1}{1-q+\epsilon/2}, \qquad q<1+\epsilon/2,\)
so that $q^{(\ell+1)}=\mathsf Q_\epsilon(\mathsf Q_\epsilon(q^{(\ell)}))$. Set
\(A_\ell\eqdef 1-q^{(\ell)}+\frac{\epsilon}{2}, \qquad A_\star\eqdef 1-q_\star+\frac{\epsilon}{2} = \frac{\epsilon+\sqrt{\epsilon^2+16}}{4}.\)
Then $q^{(\ell)}\to q_\star=1+\epsilon/2-A_\star$, and the convergence obeys the exact cross-ratio identity
\(\frac{A_\ell-A_\star}{A_\ell+A_\star^{-1}} = A_\star^{-4\ell} \frac{A_0-A_\star}{A_0+A_\star^{-1}}.\)
In particular,
\[
0\leq q_\star-q^{(\ell)}
\leq
q_\star
\left(\frac{4}{\epsilon+\sqrt{\epsilon^2+16}}\right)^{4\ell}.
\]
\end{prop}
\begin{proof}
Completing the square in $\int\exp((q y^2-(x-y)^2)/\epsilon)\d\Gaussian(0,1)(y)$ gives $\mathsf Q_\epsilon(q)$. Write $a=\epsilon/2$. Applying this map twice shows that $A_\ell=1-q^{(\ell)}+a$ satisfies
\(A_{\ell+1}=M(A_\ell), \qquad M(A)\eqdef a+\frac{A}{1+aA}.\)
The fixed points of $M$ solve $A^2-aA-1=0$ and are $A_\star>0$ and $-A_\star^{-1}$. Direct algebra for this M\"obius map gives \(\frac{M(A)-A_\star}{M(A)+A_\star^{-1}} = A_\star^{-4} \frac{A-A_\star}{A+A_\star^{-1}},\)
which proves the exact identity by iteration. Moreover, $A_0=1+a>A_\star$ and $M$ leaves $[A_\star,A_0]$ invariant. On this interval, \(0<M'(A)=\frac{1}{(1+aA)^2} \leq \frac{1}{(1+aA_\star)^2} =A_\star^{-4},\)
where $1+aA_\star=A_\star^2$. Hence $0\leq A_\ell-A_\star\leq A_\star^{-4\ell}(A_0-A_\star)$. Since $A_0-A_\star=q_\star$ and $A_\ell-A_\star=q_\star-q^{(\ell)}$, this is the announced bound.
\end{proof}
\subsection{Continuous \texorpdfstring{$\varepsilon$}{epsilon}-Sinkhorn Flow}\label{sec-continuous-epsilon-sinkhorn}
\paragraph{Parabolic Monge--Amp\`ere limit.}
Take the quadratic torus cost $c(x,y)=d_{\TT^d}(x,y)^2/2$. In the zero-temperature limit, the soft transforms concentrate near the map $x\mapsto x+\nabla u(x)$, while the first Laplace correction records its Jacobian determinant. Rescaling the accumulated correction in iteration time leads to the following flow.

\begin{defn}[Continuous \texorpdfstring{$\varepsilon$}{epsilon}-Sinkhorn flow]\label{def-continuous-epsilon-sinkhorn}
Let $\al$ and $\be$ be probability measures on the unit-volume flat torus $\TT^d$ with smooth positive densities $\density{\al}=e^{-F}$ and $\density{\be}=e^{-G}$. A gauge-fixed potential $u_t:\TT^d\to\RR$, with $\int_{\TT^d}u_t\d x=0$, follows the continuous $\epsilon$-Sinkhorn flow if
\begin{equation}\label{eq-continuous-epsilon-sinkhorn-pde}
\partial_t u_t(x)
=
\log\det\bigl(\Id+\nabla^2u_t(x)\bigr)
-G\bigl(x+\nabla u_t(x)\bigr)
+F(x)-\bar r_t,
\end{equation}
where $x+\nabla u_t(x)$ is understood modulo $\ZZ^d$, $\Id+\nabla^2u_t\succ0$, and $\bar r_t$ is the unique scalar preserving the mean-zero gauge.
\end{defn}
\begin{prop}[Stationary continuous Sinkhorn potentials]\label{prop-continuous-sinkhorn-stationary}
Assume that a smooth stationary solution of~\eqref{eq-continuous-epsilon-sinkhorn-pde} satisfies $\Id+\nabla^2u\succ0$, and define $T(x)=x+\nabla u(x)$ modulo $\ZZ^d$. Then $T_\sharp\al=\be$. Conversely, any smooth map of this form satisfying $\Id+\nabla^2u\succ0$ and $T_\sharp\al=\be$ solves the stationary equation, up to the additive gauge of $u$.
\end{prop}
\begin{proof}
At stationarity, the right-hand side before subtracting $\bar r_t$ is a constant $r$, hence
\(e^{-G(T(x))}\det\nabla T(x)=e^{-F(x)}e^r.\)
The positive Hessian condition makes $T$ an orientation-preserving local diffeomorphism of the torus, homotopic to the identity and therefore of degree one. It is thus a diffeomorphism. Integrating the displayed identity and using that both measures have unit mass gives $e^r=1$. The remaining equality is precisely the change-of-variables formula for $T_\sharp(e^{-F}\d x)=e^{-G}\d x$. Reversing the argument proves the converse.
\end{proof}
\paragraph{Rescaled continuous Sinkhorn iterates.}
Let $\mathsf S_\epsilon$ denote one complete continuous dual Sinkhorn cycle~\eqref{eq-continuous-dual-sinkhorn-iteration}, written in the sign convention $u=-f$. Using the soft transforms of Definition~\ref{def-continuous-soft-c-transform},

\(\mathsf S_\epsilon u \eqdef -\Big((-u)^{c,\epsilon}\Big)^{\bar c,\epsilon}.\)
\begin{prop}[Rescaled continuous Sinkhorn limit]\label{prop-scaled-log-sinkhorn-limit}
Let $u_\epsilon^{(0)}=u_0$ be smooth and mean zero, generate
\(u_\epsilon^{(\ell+1)} = \mathsf S_\epsilon u_\epsilon^{(\ell)} - \int_{\TT^d}\mathsf S_\epsilon u_\epsilon^{(\ell)}\d x,\)
and define $u_\epsilon(t)\eqdef u_\epsilon^{(\lfloor t/\epsilon\rfloor)}$. In particular, for $\epsilon=1/k$, one observes $u_{1/k}^{(\lfloor kt\rfloor)}$ at time $t$. Assume that, as $\epsilon\downarrow0$, these interpolants converge smoothly on compact time intervals to $u_t$, that $\Id+\nabla^2u_t\succ0$, and that the Laplace expansion below is uniform along the sequence. Then $u_t$ solves~\eqref{eq-continuous-epsilon-sinkhorn-pde}.
\end{prop}
\begin{proof}
The Laplace expansion of the two soft transforms gives the first-order consistency formula
\[
\mathsf S_\epsilon u-u
=
\epsilon\left[
\log\det(\Id+\nabla^2u)
-G(\Id+\nabla u)+F
\right]
+O(\epsilon^2),
\]
up to an $x$-independent term, which disappears under the gauge normalization. Dividing the normalized increment by $\epsilon$ and passing formally to the smooth limit gives~\eqref{eq-continuous-epsilon-sinkhorn-pde}; the subtracted spatial mean becomes $\bar r_t$.
\end{proof}
\paragraph{One-dimensional Gaussian closure.}
\(\partial_tu_t(x) = \log\bigl(1+u_t''(x)\bigr) -G\bigl(x+u_t'(x)\bigr)+F(x)-\bar r_t.\)
Although it was introduced on the torus, the same formal equation applies on $\RR$ for confining Gaussian densities. Let $\al=\Gaussian(m_\al,\sigma_\al^2)$ and $\be=\Gaussian(m_\be,\sigma_\be^2)$.

\(x+u_t'(x)=q_t+a_t(x-m_\al), \qquad a_t>0.\)
\eql{\label{eq-continuous-sinkhorn-gaussian-1d}
\dot a_t
=
\frac{1}{\sigma_\al^2}-\frac{a_t^2}{\sigma_\be^2},
\qquad
\dot q_t
=
-\frac{a_t}{\sigma_\be^2}(q_t-m_\be).
}